\section{STREAMING FEDERATED LEARNING FOR SATELLITE FAULT DIAGNOSIS}\label{sec:proposed_framework}
% Framework overview and workflow diagram

% Detailed implementation stages
The streaming federated learning framework for satellite fault diagnosis 
is illustrated in Fig. \ref{fig:streaming-federated-learning-framework}. 
Whenever a new type of fault emerges at a communication-constrained TT\&C center 
or a new communication-constrained TT\&C center joins the system, the framework is re-initiated. The convergence time of federated learning is significantly shorter than the interval between new fault occurrences, enabling continuous execution of the framework in a streaming fashion. The framework operates through six distinct stages:

\begin{figure}[t]
	\centering
	\includegraphics[width=\columnwidth]{fig/Proposed Streaming Federated Learning Framework/Streaming federated learning framework.pdf}
	\caption{Streaming federated learning framework.}
	\label{fig:streaming-federated-learning-framework}
\end{figure}

\textit{Stage 1}: Unconstrained TT\&C centers perform standard federated training~\cite{mcmahan2017communication}: each center downloads the global model, trains locally, and uploads updated parameters; the server aggregates and redistributes until convergence.

\textit{Stage 2}: The converged global model is disseminated to every communication-constrained center as the foundational knowledge base for incremental learning.

\textit{Stage 3}: Execute the IL-DDKD algorithm (Section~\ref{section:3-1}). The global model operates in inference mode, guiding communication-constrained centers to learn partial models that handle new fault types while preserving prior knowledge.

\textit{Stage 4}: Each constrained center uploads its partial model to the central server, which maintains one model per center and replaces older versions with newly uploaded ones.

\textit{Stage 5}: Execute the FL-PFC algorithm (Section~\ref{section:3-2}). Stored partial models participate in federated learning, assisting unconstrained centers with training and contributing to global model aggregation. The global model is continuously refined, transitioning from its initial to an enhanced state.

\textit{Stage 6}: The central server periodically synchronizes the enhanced global model with all constrained centers, at a rate adaptable to network conditions.

% Algorithm 1: Incremental Learning Algorithm Based on Dual-Domain Knowledge Distillation
\subsection{Incremental Learning Algorithm Based on Dual-Domain Knowledge Distillation}\label{section:3-1}
Building on Learning Without Forgetting~\cite{li2017learning} and knowledge distillation~\cite{hinton2015distilling}, we develop a dual-domain distillation technique for class-incremental satellite fault diagnosis. The teacher is the global model and the student is the partial model under training (Fig.~\ref{fig:knowledge-distillation}). To address the teacher-student capacity gap~\cite{mirzadeh2020improved}, we employ module sharing, and for new fault categories, output-layer nodes are added with randomly initialized weights.

\begin{figure}[t]
	\centering
	\includegraphics[width=\columnwidth]{fig/Proposed Streaming Federated Learning Framework/Incremental learning algorithm based on dual-domain knowledge distillation.pdf}
	\caption{Incremental learning algorithm based on dual-domain knowledge distillation.}
	\label{fig:knowledge-distillation}
\end{figure}

The loss function is:

\textbf{Supervised Loss}: This loss function is formulated to allow the partial model 
to efficiently learn fault data from the newly added TT\&C center. 
The cross entropy loss function is utilized for this objective:
\begin{equation}
	{\mathcal {L_C}}\left( {{\theta _{stu}}; {{\mathscr D}_{new}}} \right) =  - {{\bf{1}}_{{\mathcal Y_{new}}}} \log { \tilde{\mathcal Y}_{new}},
	\label{eq:LC}
\end{equation}
\noindent where ${\theta _{stu}}$ denotes the parameters of the student model,
${{\bf{1}}_{{\mathcal Y_{new}}}}$ represents the one-hot encoding of the ground truth label ${\mathcal Y_{new}}$,
${\tilde{\mathcal Y}_{new}}$ represents the model output after Softmax activation processing,
and ${{\mathscr D}_{new}}$ denotes the fault dataset from the new TT\&C center 
containing input features ${\mathcal X_{new}}$ and their corresponding true labels ${\mathcal Y_{new}}$.

\textbf{Dual domain distillation loss}: 
This loss component enables the partial model to continually assimilate fault knowledge from other TT\&C centers via the global model. 
Knowledge transfer is achieved by forwarding the attention maps produced by the intermediate layers of the partial model to the local model, 
directing the local model to concentrate on identical areas as the partial model, 
thus enabling data-domain distillation. 
Given that straightforward 1:1 alignment between feature representations of teacher and student models lacks theoretical grounding, 
we incorporate a regressor between them to enhance the alignment procedure \cite{romero2014fitnets}. 
Such architecture offers adaptability in the student model's training phase, 
preventing simple replication of the teacher model's representations by the student model.
The loss function for data-domain distillation ${\mathcal L_{f}}$ is defined as follows:
\begin{equation}
	{\mathcal L_{f}} = \frac{1}{2}{\left\| {u\left( {{{\theta _{tea}}}; {\mathcal X_{new}}} \right) - \gamma \left( {v\left( {\theta _{stu}; {\mathcal X_{new}}} \right); {W_r}} \right)} \right\|^2},
	\label{eq:LKD-1}
\end{equation}
\noindent where ${\theta _{tea}}$ represents the teacher model parameters, 
$\gamma $ denotes the regressor, ${W_r}$ refers to the parameters of the regressor, 
$u$ and $v$ represent the feature vectors generated by the teacher model and the student model, respectively. 

Meanwhile, this paper also performs distillation in the parameter domain by imposing constraints on key parameters 
to limit potential drastic changes during the learning process, 
thereby mitigating the interference of new knowledge with existing knowledge. 
The distillation loss function in the parameter domain ${\mathcal L_{par}}$ is defined as follows:
\begin{equation}
	{\mathcal L_{par}} = \frac{1}{2}\sum\limits_i {{\mathcal F_i}{{\left( {\theta _{stu}^i - \theta _{tea}^i} \right)}^2}},
	\label{eq:LKD-2}
\end{equation}
where $\mathcal F\left( \cdot \right)$ denotes the Fisher information matrix, and $i$ indicates the index of each parameter. 
The data-domain distillation loss in \eqref{eq:LKD-1} and the parameter-domain distillation loss in \eqref{eq:LKD-2}
are combined as the loss function of dual-domain distillation:
\begin{equation}
	{\mathcal {L_{KD}}} = {\mathcal L_{f}}\left( {{\theta _{stu}};{\theta _{tea}};{W_r},{\mathcal X_{new}}} \right) + {\mathcal L_{par}}\left( {{\theta _{stu}};{\theta _{tea}}} \right).
\end{equation}

The combined training objective is:
\begin{equation}
	{\mathcal L_{Dual}} = {\mathcal {L_C}}\left( {{\theta _{stu}};{{\mathscr D}_{new}}} \right) + {\lambda _o}{\mathcal {L_{KD}}}\left( {{\theta _{stu}};{\theta _{tea}};{W_r},{\mathcal X_{new}}} \right),
\end{equation}
where ${\lambda _o} > 0$ denotes the weighting parameter employed to balance the algorithm's performance between previous faults and new faults.

The weight coefficient $\lambda_o$ is adaptively adjusted: after fine-tuning the teacher to obtain a baseline accuracy $\mathcal{A}^*$, $\lambda_o$ is initialized large and decayed via $\mathcal{G}$ whenever the student's accuracy falls below $(1-\varTheta)\mathcal{A}^*$ (Algorithm~\ref{Alg:IL-DDKD}).

\begin{algorithm}[t]
	\caption{IL-DDKD Algorithm}
	\label{Alg:IL-DDKD}
	\begin{algorithmic}[1]
		\REQUIRE New TT\&C center fault data ${\mathcal X_{new}}$ and ground truth ${\mathcal Y_{new}}$, global model ${\theta _{tea}}$, local epochs $E$, local minibatch size $B$, accuracy drop margin $\Theta $, decay function $\mathcal G$, local learning rate $\eta_L $
		\ENSURE Partial model parameters $\theta _{{stu}}^ * $
		\STATE \textbf{Initialization}: Set the initial weight coefficient ${\lambda _o} = 100$, ${\mathcal A^ * } = 0$, ${W_r}$ to small random values.
		\STATE ${\mathcal A^ * } \leftarrow $Finetune$\left( {{{\mathscr D}_{new}},\eta_L ;{\theta _{tea}}} \right)$
		\WHILE {$\mathcal A < \left( {1 - \Theta } \right){\mathcal A^ * }$}
		\STATE $\mathcal B \leftarrow $ (Split ${{\mathscr D}_{new}}$ into batches of size $B$)
		\FOR {$e = 1, 2, 3,  ... , E$  }
		\FOR {batch $b \in \mathcal B$}
		\STATE ${\mathcal {L_C}} \leftarrow $CE$ ( {{{\bf{1}}_{{\mathcal Y_{new}}}},{\mathcal{\tilde Y}_{new}}} )$
		\STATE $teacher\_out \leftarrow u\left( {{\theta _{tea}},{\mathcal X_{new}}} \right)$
		\STATE $student\_out \leftarrow \gamma \left( {v\left( {{\theta _{stu}},{\mathcal X_{new}}} \right),{W_r}} \right)$
		\STATE ${L_{f}} \leftarrow$ MSE $\left( {teacher\_out, student\_out} \right)$
		\STATE ${\mathcal L_{par}} \leftarrow \mathcal F\left( {{\theta _{stu}},{\theta _{tea}}} \right)$
		\STATE ${\mathcal{L_{KD}}} = {\mathcal L_{f}} + {\mathcal L_{par}}$
		\STATE ${\mathcal L_{Dual}} = {\mathcal {L_C}} + {\lambda _o}{\mathcal {L_{KD}}}$
		\STATE ${\theta _{stu}} \leftarrow {\theta _{stu}} - \eta_L \nabla {\mathcal L_{Dual}}$
		\ENDFOR
		\ENDFOR
		\STATE $\mathcal A \leftarrow $\textbf{Evaluate}$\left( {{{{\mathscr D}_{new}}};{\theta _{stu}}} \right)$
		\IF {$\mathcal A < \left( {1 - \Theta } \right){\mathcal A^ * }$}
		\STATE ${\lambda _o} \leftarrow \mathcal G\left( {{\lambda _o}} \right)$
		\ENDIF
		\ENDWHILE
		\STATE $\theta _{{{stu}}}^ *  \leftarrow {\theta _{stu}}$
		\RETURN $\theta _{{{stu}}}^ * $
	\end{algorithmic}
\end{algorithm}

% Algorithm 2: Federated Learning Algorithm Based on Partial Forgetting Compensation
\subsection{Federated Learning Algorithm Based on Partial Forgetting Compensation}\label{section:3-2}
\begin{figure}[t]
	\centering
	\includegraphics[width=\columnwidth]{fig/Proposed Streaming Federated Learning Framework/Federated learning algorithm based on local forgetting compensation.pdf}
	\caption{Federated learning algorithm based on partial forgetting compensation.}
	\label{fig:partial-forgetting-compensation}
\end{figure}

To enable the knowledge acquired at communication-constrained TT\&C centers 
to be utilized in subsequent federated learning processes of unconstrained TT\&C centers, 
this paper implements a virtual client as a proxy server for each communication-constrained TT\&C center. 
The proxy server stores the model trained by the communication-constrained TT\&C center 
and applies partial forgetting compensation to the entire system, 
ensuring that model parameter updates remain within the target range.
For privacy (see Section~\ref{sec:Related_Work}), the proxy server grants white-box access only to the central server and black-box access to individual centers. Fig.~\ref{fig:partial-forgetting-compensation} illustrates the architecture and Algorithm~\ref{Alg:FL-PFC} details the execution. Each round consists of four steps:

\textbf{Global Model Download}: The server randomly selects a subset of centers and distributes the global model $\omega_t$, following the standard partial-participation scheme~\cite{mcmahan2017communication}.

\textbf{Local Model Update}: The selected TT\&C centers utilize ${\omega _t}$ as the initial model and perform training 
on their respective local fault data ${{\mathscr D}_{i}}$ using the \textit{inheritance distillation method} 
described in Section~\ref{sec:proposed_framework}-\ref{section:3-2}-\ref{section:3-2-1}, generating an updated local model $\omega _{t + 1}^i$.

\textbf{Local Model Upload}: Participating TT\&C centers transmit their updated local model parameters to the central server.

\textbf{Global Model Update}: The central server consolidates model parameters received from all participating TT\&C centers 
and refreshes the global model employing the \textit{Dynamic Weighting Method} elaborated in Section~\ref{sec:proposed_framework}-\ref{section:3-2}-\ref{section:3-2-2}.

\begin{algorithm}[t]
	\caption{FL-PFC Algorithm}
	\label{Alg:FL-PFC}
	\begin{algorithmic}[1]
		\REQUIRE The $i$-th TT\&C center fault data ${{\mathscr D}_{i}}$, the number of TT\&C centers $K$, communication rounds $T$, local epochs $E$, fraction of TT\&C centers on each round $p$, model saved in proxy server $\theta _ \times ^ * $, global learning rate $\eta $
		\ENSURE Global model parameters ${\omega ^ * }$
		\FOR {$t = 1, 2, 3,  ... , T$}
		\STATE $m \leftarrow \max \left( {p \cdot K,1} \right)$
		\STATE ${S_t} \leftarrow $(randomly selected subset of $m$ TT\&C centers)
		\FOR {$i \in {S_t}$ \textbf{in parallel}}
		\STATE $\omega _{t+1}^i \leftarrow $\textbf{Inheritance Distillation Method}(${{\mathscr D}_{i}}$; $\theta _ \times ^ * $; ${\eta}$; $E$)
		\ENDFOR
		\STATE ${\omega _{t + 1}} \leftarrow $\textbf{Dynamic Weighting Method}(${{\mathscr D}_{i}}$;${S_t}$; $\omega _{t}$; $\omega _{t+1}^i$)
		\ENDFOR
		\STATE ${\omega ^ * } \leftarrow {\omega _T}$
		\RETURN ${\omega ^ * }$
	\end{algorithmic}
\end{algorithm}

% Inheritance Distillation Method
\subsubsection{Inheritance Distillation Method}
\label{section:3-2-1}
This section maintains the teacher-student model paradigm. 
However, during this stage, the roles of teacher and student models are reversed: 
the teacher model signifies the model developed by the communication-limited TT\&C center,
whereas the student model transforms into the model developed by unconstrained TT\&C centers. 
As the entire mechanism embodies the progressive dissemination and preservation of knowledge, 
we designate it as ``\textbf{inheritance distillation}". 
Within this configuration, the teacher model aligns with the trained partial model (blue) illustrated in Fig.~\ref{fig:streaming-federated-learning-framework}, 
and the student model matches the original global model (gray) depicted in Fig.~\ref{fig:streaming-federated-learning-framework}.

During the inheritance distillation process, the teacher model is managed by the proxy server, 
which assists each unconstrained TT\&C center in processing local fault data and training the model in a black-box manner. 
The purpose is to separate the model from the knowledge representation, avoiding the exposure of sensitive information and achieving privacy protection. 
The algorithm's loss function is formulated as follows:

\textbf{Classification Loss}: This loss function is designed to enable the local model to effectively learn the fault data ${{\mathscr D}_i}$ from its respective TT\&C center. 
The classification loss $\mathcal L_C^i$ for TT\&C center $i$ is defined as follows:
\begin{equation}
	\label{eq:LCi}
	\mathcal L_C^i = {{\mathcal D}_{\rm{CE}}}\left( {{\mathcal P_i}\left( {\omega _{t,e}^i,{\mathcal X_i}} \right),{\mathcal Y_i}} \right),
\end{equation}
\noindent where $\omega _{t,e}^i$ represents the model parameters of TT\&C center $i$ at iteration $e$ of round $t$, 
$e$ denotes the local training iterations, 
$t$ represents the federated learning rounds, 
${\mathcal P_i}\left( {\omega _{t,e}^i, {\mathcal X_i}} \right)$ indicates the probability values predicted by the model $\omega _{t,e}^i$ for input data ${\mathcal X_i}$, 
${\mathcal Y_i}$ represents the ground truth labels, 
and ${{\mathcal D}_{\rm{CE}}}\left( { \cdot ,\cdot } \right)$ denotes the cross-entropy loss function.

\textbf{Distillation loss}: As in IL-DDKD, knowledge transfer uses intermediate-layer attention maps to align the local model's focus with the partial model. The distillation loss $\mathcal L_{\mathcal KD}^i$ is:
\begin{equation}
	\label{eq:LKDi}
	\mathcal L_{\mathcal KD}^i = \frac{1}{2}{\left\| {{u_i}\left( {\theta _ \times ^ * ,{\mathcal  X_i}} \right) - {\gamma _i}\left( {{v_i}\left( {\omega _{t,e}^i,{\mathcal  X_i}} \right),W_r^i} \right)} \right\|^2},
\end{equation}
\noindent where $\theta _ \times ^ * $ represents the model saved in proxy server, 
$\gamma_i$ is the mapper, and $W^i_r$ are the model parameters of the regressor,
${u_i}$ and ${v_i}$ are the feature vectors output by the partial model and the local model, respectively. 

The integration of the classification loss in \eqref{eq:LCi} and the distillation loss in \eqref{eq:LKDi} constitutes the complete training objective:
\begin{equation}
	{\mathcal L^i} = \mathcal L_C^i\left( {\omega _{t,e}^i,{{\mathscr D}_i}} \right) + {\lambda ^i}\mathcal L_{\mathcal KD}^i\left( {\omega _{t,e}^i;\theta _ \times ^ * ;W_r^i,{\mathcal X_i}} \right),
\end{equation}
where $\lambda^i > 0$ is the weight coefficient 
used to balance the learning performance of the local model on its own TT\&C center's fault data 
and the communication-constrained TT\&C center's fault data.

Drawing on the approach from \cite{li2020federated}, this paper introduces a regularization term into the objective function 
to mitigate the impact of data heterogeneity among TT\&C centers. The final objective function of the algorithm is expressed as follows:
\begin{equation}
	{{\tilde {\mathcal  L^i}}} = \mathcal L_C^i + {\lambda ^i}\mathcal L_{\mathcal KD}^i + \frac{\mu }{2}{\left\| {\omega _{t,e}^i - {\omega _t}} \right\|^2},
\end{equation}
where ${\omega _t}$ denotes the global model parameters disseminated by the central server during the $t$-th round, 
and $\mu$ is the penalty constant that is adaptively adjusted \cite{li2020federated}. 
Specifically, this paper adopts a simple heuristic strategy: 
when the objective function value exhibits an upward trend, the penalty constant $\mu$ is increased to enhance the constraint effect; 
conversely, when the objective function value continues to decrease, 
the penalty constant $\mu$ is appropriately reduced to accelerate convergence. 
Algorithm \ref{Alg:Inheritance_Distillation_Method} presents the complete execution process of the inheritance distillation method.

\begin{algorithm}[t]
	\caption{Inheritance Distillation Method}
	\label{Alg:Inheritance_Distillation_Method}
	\begin{algorithmic}[1]
		\REQUIRE Fault data ${{\mathscr D}_{i}}$ from the $i$-th TT\&C center, the local model achieved by the $i$-th TT\&C center in the $t$-th round $\omega _{t + 1}^i$, global model for the $t$-th round ${\omega _t}$, the collection of TT\&C centers randomly chosen in the $t$-th round ${S_t}$
		\ENSURE The local model parameters update of the $i$-th TT\&C center $\omega _{t + 1}^i$
		\STATE \textbf{Initialization}: Set the weight coefficient ${\lambda ^i} = 1$, $W_r^i$ to small random values, $\omega _{t,0}^i = {\omega _t}$.
		\STATE $\mathcal B \leftarrow $ (Split ${{\mathscr D}_{i}}$ into batches of size $B$)
		\FOR {$e = 1, 2, 3,  ... , E$}
		\FOR {batch $b \in \mathcal B$}
		\STATE $\mathcal  L_C^i \leftarrow$ CE$\left( {{\mathcal P_i}\left( {\omega _{t,e}^i,{\mathcal X_i}} \right),{\mathcal Y_i}} \right)$
		\STATE ${\mathcal L_{\mathcal KD}^i} \leftarrow$ MSE$\left( {{u_i}\left( {\theta _ \times ^ * ,{\mathcal X_i}} \right),{\gamma _i}\left( {{v_i}\left( {\omega _{t,e}^i,{\mathcal X_i}} \right),W_r^i} \right)} \right)$
		\STATE ${ {\tilde { \mathcal L}^i}} = \mathcal L_C^i + {\lambda ^i} {\mathcal L_{\mathcal KD}^i} + \frac{\mu }{2}{\left\| {\omega _{t,e}^i - {\omega _t}} \right\|^2}$
		\STATE $\omega _{t,e + 1}^i \leftarrow \omega _{t,e}^i - {\eta}\nabla  {\tilde {\mathcal L}^i}$
		\ENDFOR
		\ENDFOR
		\STATE $\omega _{t + 1}^i \leftarrow \omega _{t,E}^i$
		\RETURN $\omega _{t + 1}^i$
	\end{algorithmic}
\end{algorithm}

% Dynamic Weighting Method
\subsubsection{Dynamic Weighting Method}
\label{section:3-2-2}
Following each client $k$'s local model update, 
the central server collects these model parameters and 
executes a weighted average to update the global model ${\omega _{t + 1}} \leftarrow \sum\nolimits_{k = 1}^K {\frac{{{n_k}}}{n}\omega _{t + 1}^k} $, 
which represents the standard approach recommended by traditional federated learning algorithms \cite{von2018data}. 
However, this conventional aggregation method demonstrates poor performance in our study, 
primarily because it assigns weights to each client based exclusively on data volume, 
without adequately considering the importance and individual characteristics of each client during global model aggregation. 
The distinctive challenge in our research stems from the data heterogeneity of data across TT\&C centers, 
which results in unsatisfactory global model performance.

\begin{algorithm}[t]
	\caption{Dynamic Weighting Method}
	\label{Alg:Dynamic_Weighting_Method}
	\begin{algorithmic}[1]
		\REQUIRE The $i$-th TT\&C center fault data ${{\mathscr D}_{i}}$, the local model obtained by the $i$-th TT\&C center in the $t$-th round $\omega _{t + 1}^i$, global model of the $t$-th round ${\omega _t}$, the set of TT\&C centers randomly selected in the $t$-th round ${S_t}$
		\ENSURE The global model parameters update of the central server ${\omega _{t + 1}}$
		\STATE $u_{t + 1}^i \leftarrow $Test (${{\mathscr D}_{i}}$, $\omega _{t + 1}^i$)
		\STATE $v_t^i \leftarrow $ Test (${{\mathscr D}_{i}}$, ${\omega _t}$)
		\STATE ${\gamma ^i} \leftarrow \% \; ({{e^{u_{t + 1}^i - v_t^i}}})  $
		\STATE ${\pi _{proxy}},{\gamma _t^{proxy}} \leftarrow \bar \pi ,\bar C_{\gamma}(t+1)^{-(1+\rho )}$
		\STATE ${\omega _{t + 1}} \leftarrow \sum\nolimits_{i \in {S_t} \cup proxy} {\phi \left[\kern-0.15em\left[ {{\pi _i}{\gamma ^i}}  \right]\kern-0.15em\right]\omega _{t + 1}^i} $
		\RETURN ${\omega _{t + 1}}$
	\end{algorithmic}
\end{algorithm}

This paper introduces an importance balancing factor $\gamma$ (Algorithm~\ref{Alg:Dynamic_Weighting_Method}) in the federated learning aggregation process 
to enable the global model to dynamically adjust the importance of knowledge contributions from different TT\&C centers throughout the learning process. 
The importance balancing factor $\gamma$ is defined as:
\begin{equation}
	{\gamma ^k} = \frac{{{e^{u_{t + 1}^k - v_t^k}}}}{{\sum\nolimits_{k \in {S_t}} {{e^{u_{t + 1}^k - v_t^k}}} }},
\end{equation}
\noindent where $u_{t + 1}^k$ and $v_t^k$ represent the classification accuracy of the local model at round $t + 1$ and the global model at round $t$ on the fault dataset ${{\mathscr D}_{k}}$ of the $k$-th TT\&C center, respectively. The global model parameter update on the central server is computed according to the following equation:
\begin{equation}
	{\omega _{t + 1}} \leftarrow \sum_{k \in {S_t} \cup proxy} {\phi \left[\kern-0.15em\left[ {{\pi _k}{\gamma ^k}} 
		\right]\kern-0.15em\right]\omega _{t + 1}^k} ,
\end{equation}
\noindent where ${\pi _k}$ denotes the proportion of fault data in each TT\&C center, 
and $\phi \left[\kern-0.15em\left[ \cdot \right]\kern-0.15em\right]$ represents the normalized mapping in the metric space.
Interestingly, in the model aggregation process, 
local models with poorer performance are intentionally assigned higher weights. 
An intuitive explanation is that models requiring more learning in previous rounds 
should receive greater attention in subsequent training iterations. 
This design effectively compensates for weak performers 
and thereby enhances the overall performance and robustness of the model.

Furthermore, throughout the entire federated learning process, 
unconstrained TT\&C centers gradually acquire 
the fault knowledge of the communication-constrained TT\&C center with assistance from the virtual client. 
As these centers continuously utilize TT\&C data from unconstrained centers for training, 
they can further improve their performance on local fault datasets. 
Consequently, the performance of local models trained by unconstrained TT\&C centers 
in each round will progressively surpass the model performance of the virtual client. 
At this stage, the virtual client essentially becomes a performance bottleneck. 
To address this issue, this paper incorporates an explicit polynomial decay schedule
\begin{equation}
	\gamma _t^{proxy}=\bar C_{\gamma}(t+1)^{-(1+\rho )}, \qquad \rho > 0,
\end{equation}
into the weight coefficient of the communication-constrained TT\&C center 
to mitigate the negative impact of the virtual client during model aggregation. This explicit schedule is also the form required by the convergence analysis to make the proxy residual term summable.